% A minimizing sequence xk=1/k for phi=x^2 on S=[0,1] illustrates the two proof links.
\begin{tikzpicture}[font=\small,>={Latex[length=2mm]}]
 \node at (2,2) {compact domain $S=[0,1]$};
 \draw[very thick,black] (0,0)--(4,0);
 \fill (0,0) circle(2pt) node[below] {$\hat x=0\in S$};
 \fill (4,0) circle(2pt) node[below] {$1$};
 \foreach \x in {2,1.3333,1,.8,.6667,.5,.4}{\fill[blue!70!black] (\x,0) circle(1.6pt);}
 \draw[->,thick] (2.4,.6)--(.2,.6) node[midway,above] {$x^{k_i}\to\hat x$};
 \node[align=center] at (2,-1) {compactness retains the limit\\inside the domain};
 \draw[->,thick] (4.9,.5)--(6.6,.5) node[midway,above] {$\phi(x)=x^2$};
 \begin{scope}[xshift=8cm]
  \node at (0,2) {function values};
  \draw[->,black] (0,-.3)--(0,1.6);
  \fill (0,0) circle(2pt) node[right] {$\phi(\hat x)=0$};
  \foreach \y in {1,.25,.111111,.0625,.04,.027778,.015625,.01}{\fill[blue!70!black] (0,\y) circle(1.6pt);}
  \draw[->,thick] (-.7,1.45)--(-.7,.12);
  \node[align=left,anchor=west] at (.5,.9) {$\phi(x^{k_i})\to m$\\and $\phi(x^{k_i})\to\phi(\hat x)$};
  \node[align=center] at (.5,-1) {continuity preserves\\the limit of function values};
 \end{scope}
 \node at (4.9,-1.85) {the minimizing sequence has $x^k=1/k$, $\phi(x^k)=1/k^2\to m=0$};
\end{tikzpicture}
